\chapter[Lezione 13 --- 23 ottobre 2026]{Lezione 13: funzione di Dirichlet e proprietà dei limiti}
\label{chap:lezione-13}

\section{Funzione di Dirichlet}
\label{sec:lez13-dirichlet}

\begin{example}[Funzione di Dirichlet]
\label{ex:lez13-dirichlet}
Si consideri la funzione \(f\colon\R\longrightarrow\R\) definita da
\[
  f(x):=
  \begin{cases}
    1, & x\in\R\setminus\Q,\\
    -1, & x\in\Q.
  \end{cases}
\]
Si ha \(f(0)=-1\), ma \(f\) non ammette limite per \(x\to0\).
\end{example}

\begin{proof}
Il dominio è \(A=\R\). Per ogni \(\delta>0\), il punto
\(a=\delta/2\) appartiene ad \(A\) e soddisfa
\[
  0<|a|=\frac{\delta}{2}<\delta.
\]
Pertanto \(0\) è punto di accumulazione per \(A\).
Si supponga anzitutto che
\[
  \lim_{x\to0}f(x)=\ell\in\R.
\]
Per la \cref{def:lez08-limite}, scegliendo
\(\varepsilon=1/2\), esiste \(\delta>0\) tale che
\[
  |f(x)-\ell|<\frac12
  \qquad\forall x\in\R:\ 0<|x|<\delta.
\]
Per il \cref{thm:razionali-irrazionali-tra-reali}, esistono
\(r\in\Q\) e \(s\in\R\setminus\Q\) tali che
\(0<r,s<\delta\). Poiché \(f(r)=-1\) e \(f(s)=1\), si ottiene
\[
  |-1-\ell|<\frac12,
  \qquad
  |1-\ell|<\frac12.
\]
Per la disuguaglianza triangolare, ne segue
\[
  2=|(1-\ell)-(-1-\ell)|
  \le |1-\ell|+|-1-\ell|<1,
\]
una contraddizione.

Neppure \(+\infty\) e \(-\infty\) possono essere limiti:
poiché \(-1\le f(x)\le1\) per ogni \(x\in\R\),
la \cref{def:lez08-limite} non è soddisfatta scegliendo
rispettivamente la soglia \(K=2\) o \(K=-2\).
Pertanto \(f\) non ammette limite per \(x\to0\).
\end{proof}

\section{Proprietà elementari dei limiti}
\label{sec:lez13-proprieta-limiti}

\begin{proposition}[Proprietà elementari dei limiti]
\label{prop:lez13-proprieta-limiti}
Sia \(f\colon A\longrightarrow\R\), con \(A\subseteq\R\)
non vuoto, e sia \(x_0\in\overline{\R}\) un punto di
accumulazione per \(A\). Allora:
\begin{enumerate}
  \item per ogni \(\ell\in\R\),
  \[
    \lim_{x\to x_0}f(x)=\ell
    \quad\Longleftrightarrow\quad
    \lim_{x\to x_0}(f(x)-\ell)=0;
  \]
  \item
  \[
    \lim_{x\to x_0}f(x)=0
    \quad\Longleftrightarrow\quad
    \lim_{x\to x_0}|f(x)|=0;
  \]
  \item per ogni \(\ell\in\R\setminus\{0\}\),
  \[
    \lim_{x\to x_0}f(x)=\ell
    \quad\Longrightarrow\quad
    \lim_{x\to x_0}|f(x)|=|\ell|;
  \]
  \item
  \[
    \lim_{x\to x_0}f(x)=0
    \quad\Longleftrightarrow\quad
    \lim_{x\to x_0}f(x)^2=0.
  \]
\end{enumerate}
\end{proposition}

\begin{proof}
\emph{Primo punto.}
Per ogni \(x\in A\) e \(\ell\in\R\),
\[
  |(f(x)-\ell)-0|=|f(x)-\ell|.
\]
Le due condizioni della \cref{def:lez08-limite}
coincidono, quindi vale l'equivalenza.

\emph{Secondo punto.}
Per ogni \(x\in A\),
\[
  \bigl||f(x)|-0\bigr|=|f(x)-0|.
\]
Anche in questo caso le condizioni della
\cref{def:lez08-limite} coincidono.

\emph{Terzo punto.}
Si supponga \(f(x)\to\ell\in\R\setminus\{0\}\)
per \(x\to x_0\). Per la disuguaglianza triangolare inversa,
\[
  \bigl||f(x)|-|\ell|\bigr|\le|f(x)-\ell|.
\]
Dato \(\varepsilon>0\), per la \cref{def:lez08-limite}
esiste un intorno \(U_{x_0}\) tale che
\[
  |f(x)-\ell|<\varepsilon
  \qquad\forall x\in(U_{x_0}\cap A)\setminus\{x_0\}.
\]
Per gli stessi \(x\) si ha
\[
  \bigl||f(x)|-|\ell|\bigr|<\varepsilon,
\]
e quindi \(\lim_{x\to x_0}|f(x)|=|\ell|\).

\emph{Quarto punto.}
Si dimostrano separatamente le due implicazioni.

\emph{Prima implicazione \((\Rightarrow)\).}
Si supponga \(f(x)\to0\) per \(x\to x_0\).
Per ogni \(\varepsilon>0\), applicando la
\cref{def:lez08-limite} con \(\sqrt{\varepsilon}>0\),
esiste un intorno \(U_{x_0}\) tale che
\[
  |f(x)|<\sqrt{\varepsilon}
  \qquad\forall x\in(U_{x_0}\cap A)\setminus\{x_0\}.
\]
Ne segue
\[
  |f(x)^2|=|f(x)|^2<\varepsilon
  \qquad\forall x\in(U_{x_0}\cap A)\setminus\{x_0\}.
\]
Pertanto \(\lim_{x\to x_0}f(x)^2=0\).

\emph{Seconda implicazione \((\Leftarrow)\).}
Si supponga \(f(x)^2\to0\) per \(x\to x_0\).
Per ogni \(\varepsilon>0\), applicando la
\cref{def:lez08-limite} con \(\varepsilon^2>0\),
esiste un intorno \(U_{x_0}\) tale che
\[
  |f(x)^2|<\varepsilon^2
  \qquad\forall x\in(U_{x_0}\cap A)\setminus\{x_0\}.
\]
Poiché \(|f(x)^2|=|f(x)|^2\), si ottiene
\[
  |f(x)|<\varepsilon
  \qquad\forall x\in(U_{x_0}\cap A)\setminus\{x_0\}.
\]
Pertanto \(\lim_{x\to x_0}f(x)=0\).
\end{proof}
